\documentclass[11pt,a4paper,oneside]{article}
\usepackage[utf8]{inputenc}
\usepackage[italian]{babel}
\usepackage{hyperref}
\usepackage{listings}
\usepackage{xcolor}
\usepackage{geometry}
\usepackage{fancyhdr}
\usepackage{tcolorbox}
\usepackage{array}

\geometry{margin=2cm, headheight=15pt}

% Stili per il codice
\definecolor{codegreen}{rgb}{0,0.6,0}
\definecolor{codegray}{rgb}{0.5,0.5,0.5}
\definecolor{codepurple}{rgb}{0.58,0,0.82}
\definecolor{backcolour}{rgb}{0.96,0.96,0.96}

\lstdefinestyle{mystyle}{
    backgroundcolor=\color{backcolour},   
    commentstyle=\color{codegreen},
    keywordstyle=\color{magenta},
    numberstyle=\tiny\color{codegray},
    stringstyle=\color{codepurple},
    basicstyle=\ttfamily\scriptsize,
    breakatwhitespace=false,         
    breaklines=true,                 
    captionpos=b,                    
    keepspaces=true,                 
    numbers=left,                    
    numbersep=5pt,                  
    showspaces=false,                
    showstringspaces=false,
    showtabs=false,                  
    tabsize=2
}
\lstset{style=mystyle}

% Intestazioni e piè di pagina
\pagestyle{fancy}
\fancyhf{}
\fancyhead[L]{Proposta Tecnica - Verifica Matricole}
\fancyhead[R]{\thepage}
\fancyfoot[C]{Associazione Culturale Gulliver - Riservato}

\title{
    \vspace{1cm}
    \textbf{\huge Proposta di Integrazione:\\ Sistema di Verifica OTP per Accesso Gruppi Matricole}\\
    \vspace{0.5cm}
    \large Analisi Tecnica, Sicurezza, Utilizzo di Redis e Aspetti Regolamentari GDPR\\
    \vspace{1cm}
}
\author{\textbf{Associazione Culturale Gulliver -- Lista di Rappresentanza UNIVPM}}
\date{\today}

\begin{document}

\maketitle

\section{Introduzione e Contesto}
L'Associazione Culturale Gulliver organizza e gestisce i gruppi WhatsApp di riferimento per tutte le matricole dell'Università Politecnica delle Marche (UNIVPM). Questa attività, volta ad agevolare l'integrazione e la comunicazione tra gli studenti, ha registrato criticità dovute all'accesso di utenti esterni all'Ateneo o malintenzionati. 

Per mitigare questo fenomeno, si è valutato l'inserimento di filtri all'ingresso chiedendo nome, cognome e numero di matricola in fase di richiesta d'accesso, per poi validare i dati manualmente tramite ricerca sulla piattaforma Microsoft Teams di Ateneo. Tuttavia, tale procedura risulta complessa, inefficiente sotto il profilo temporale e solleva perplessità di conformità legale inerenti all'acquisizione manuale di dati identificativi. 

Si propone qui una soluzione \textbf{completamente automatizzata}, basata sull'invio di un codice transitorio OTP (One-Time Password) alla casella di posta elettronica istituzionale dello studente.

\section{Architettura e Flusso di Funzionamento}
Il sistema di verifica è strutturato per operare in ambiente serverless su Vercel, sfruttando le Next.js API Routes e una base dati Redis effimera per mantenere lo stato dell'OTP.

Il flusso di interazione si articola nelle seguenti fasi:
\begin{enumerate}
    \item \textbf{Selezione del Corso e Inserimento Matricola}: Lo studente accede alla pagina \texttt{/matricole/gruppi}, seleziona il proprio corso di studi da un menu a tendina ed inserisce la propria matricola (es. \texttt{1114896}).
    \item \textbf{Generazione OTP ed Invio Email (Server-Side)}:
    \begin{itemize}
        \item Il server riceve la richiesta e valida la struttura della matricola (es. formato numerico a 7 cifre).
        \item Viene composto l'indirizzo email istituzionale dello studente secondo la regola dell'Ateneo: \texttt{s[matricola]@studenti.univpm.it} (es. \texttt{s1114896@studenti.univpm.it}).
        \item Il server genera un codice OTP casuale e sicuro a 6 cifre e lo salva temporaneamente su Redis (con una scadenza di 10 minuti).
        \item Tramite un servizio di invio email transazionale (es. API di Resend o SMTP dell'associazione), viene spedito il codice OTP all'indirizzo istituzionale.
    \end{itemize}
    \item \textbf{Inserimento e Verifica dell'OTP}: Lo studente visualizza a schermo un campo di testo per l'inserimento dell'OTP. All'invio del codice, il server confronta l'input con il valore memorizzato su Redis:
    \begin{itemize}
        \item Se l'OTP corrisponde, la chiave viene immediatamente rimossa da Redis (per impedirne il riutilizzo) e viene restituito e visualizzato a schermo il link d'accesso al gruppo WhatsApp desiderato.
        \item Se l'OTP è errato, viene incrementato un contatore di tentativi per bloccare il brute-forcing (max 3 tentativi consentiti, dopodiché l'OTP viene invalidato).
    \end{itemize}
\end{enumerate}

\section{Utilizzo di Upstash Redis}
L'infrastruttura di Gulliver Web utilizza attualmente \textbf{Upstash Redis} come database Serverless chiave-valore. Le implementazioni attuali includono:
\begin{itemize}
    \item La persistenza delle configurazioni dei form dinamici sotto la chiave \texttt{gulliver:forms}.
    \item La gestione dei popup dinamici per le elezioni studentesche, tracciando lo stato attivo, il titolo, il testo e il timestamp di versione (\texttt{popupVersion}) per resettare lo stato di visualizzazione locale degli utenti.
\end{itemize}

Nel nuovo sistema di verifica, Redis verrebbe utilizzato per due scopi fondamentali:
\subsection{Conservazione dell'OTP}
Per ogni richiesta di verifica, Redis memorizza un record transitorio del tipo:
\begin{lstlisting}[language=json]
// Chiave: otp:1114896
{
  "otp": "489201",
  "corso": "ingegneria_informatica",
  "attempts": 0
}
\end{lstlisting}
Questo record viene creato con un'opzione di scadenza automatica (\texttt{EX} o \texttt{PX} in comandi Redis) impostata a \textbf{600 secondi} (10 minuti).

\subsection{Rate Limiting e Prevenzione degli Abusi}
Per prevenire lo spam di email a indirizzi di studenti terzi o l'esaurimento dei piani di invio email, Redis tiene traccia delle soglie di accesso:
\begin{itemize}
    \item \textbf{Rate limit per IP}: Limita il numero di richieste provenienti dallo stesso indirizzo IP (es. massimo 5 richieste ogni 15 minuti).
    \item \textbf{Rate limit per Matricola}: Impedisce la generazione continua di codici OTP per la stessa matricola (es. massimo 2 tentativi di invio ogni 10 minuti).
\end{itemize}

\section{Analisi Legale e Conformità GDPR}
Il punto critico sollevato nel dibattito interno riguarda la legittimità della richiesta del numero di matricola da parte dei rappresentanti degli studenti. Di seguito si analizzano gli aspetti legali alla luce del Regolamento Generale sulla Protezione dei Dati (GDPR - Regolamento UE 2016/679).

\subsection{Matricola: Dato Personale vs Dato Sensibile}
\begin{tcolorbox}[colback=red!5!white,colframe=red!75!black,title=Distinzione Legale Fondamentale]
    Ai sensi del GDPR, il numero di matricola \textbf{non è un dato sensibile} (formalmente definito all'Art. 9 del GDPR come "categorie particolari di dati"), bensì un \textbf{dato personale comune} (Art. 4 GDPR).
\end{tcolorbox}

\begin{itemize}
    \item \textbf{Dato Personale (Art. 4 GDPR)}: È "qualsiasi informazione riguardante una persona fisica identificata o identificabile". La matricola permette l'identificazione indiretta dell'utente (attraverso l'associazione con i database dell'università), pertanto rientra in questa categoria.
    \item \textbf{Dati Sensibili (Art. 9 GDPR)}: Sono dati relativi all'origine razziale o etnica, opinioni politiche, convinzioni religiose, dati biometrici, genetici o relativi alla salute. La matricola scolastica o universitaria non contiene tali informazioni ed è quindi esclusa da questa classificazione a tutela rinforzata.
\end{itemize}

\subsection{Minimizzazione e Conservazione dei Dati (Art. 5 GDPR)}
Il sistema proposto aderisce in modo nativo ai principi di \textbf{Privacy by Design} e \textbf{Data Minimizzazione}:
\begin{itemize}
    \item \textbf{Nessun database persistente}: Il sito web non crea e non mantiene un registro degli studenti. La matricola viene memorizzata in modo volatile su Redis e scompare definitivamente entro 10 minuti.
    \item \textbf{Assenza di associazione identitaria}: Il sistema non richiede nome, cognome o altre informazioni identificative dirette. Agli occhi dell'Associazione, la matricola inserita rimane un dato pseudonimo non associato a un'identità anagrafica visibile.
\end{itemize}

\subsection{Base Giuridica del Trattamento (Art. 6 GDPR)}
Il trattamento transitorio della matricola per finalità di verifica trova fondamento nell'\textbf{Articolo 6, paragrafo 1, lettera f) del GDPR}: il \textbf{legittimo interesse} del titolare del trattamento (l'Associazione Gulliver) nel garantire la sicurezza, l'integrità e la riservatezza dei propri canali di comunicazione per gli studenti, escludendo profili non autorizzati o spammer. Questo interesse prevale sui diritti dell'interessato, dato che l'impatto sulla privacy di quest'ultimo è nullo (grazie all'immediata distruzione dei dati).

\section{Conclusioni sulla Fattibilità}
Il sistema descritto coniuga elevata usabilità per le matricole, un'efficace barriera contro le intrusioni nei gruppi ed una compliance totale alle normative europee sulla riservatezza dei dati. Grazie alla natura stateless di Next.js e all'uso combinato di Upstash Redis per la gestione di OTP effimeri e rate limiting, l'infrastruttura serverless su Vercel è ideale per ospitare questa architettura senza costi aggiuntivi.

\end{document}
